\section{Discussion}
\label{sec:discussion}

\para{What do the results mean?} The central result is not subtle. On impossible structured prompts, both tested models switch from often abstaining to always attempting an answer when abstention is forbidden. That is the cleanest evidence for the ``answer anyway'' loophole. The model is not merely making a hard factual mistake; it is choosing compliant-looking output in a setting where the request has no valid completion.

\para{Why are the factual results weaker?} The factual benchmarks reveal that answer pressure interacts with model-specific policy. \gptfivemini treats abstention instructions as meaningful guidance. When allowed to abstain, it frequently does so; when pushed to answer, it mostly stops abstaining and becomes much less reliable. \gptfourone behaves differently. It already tends to answer under all conditions, so \forcedanswer has little additional effect on factual QA. In that sense, \gptfourone exhibits a different version of the same problem: the model often answers anyway even when it is explicitly told not to guess.

\para{Relation to prior work.} These findings connect several benchmark traditions. They are consistent with the abstention failures emphasized by AbstentionBench \citep{kirichenko2025abstentionbench}, but they add a controlled manipulation that isolates answer pressure as a likely driver. They also align with hallucination work such as \halogen and HaluEval \citep{ravichander2025halogen,li2023halueval}, while clarifying that some fabricated outputs are best understood as compliance failures rather than generic knowledge errors.

\para{Limitations.} This study tests only two generation models from a single provider. The factual experiments use 50-item slices rather than full benchmark runs, and \simpleqa/\truthfulqa grading relies on an LLM judge rather than purely symbolic scoring. Prompt wording is another limitation: different abstention or pressure phrasings could change the exact rates. Finally, our definition of loophole use is behavioral and operational. It does not imply strategic intent of the kind studied in Sleeper Agents \citep{hubinger2024sleeper}.

\para{Broader implications.} The practical lesson is that systems should not reward ``doing something'' when abstention is correct. If downstream pipelines treat any answer as progress, they may systematically incentivize fabricated compliance. External validators, benchmark-specific abstention policies, and interfaces that preserve explicit uncertainty are likely to be more robust than prompt-only pressure to answer.
